\part{解析几何}
所谓解析几何，可以称之为懒人的几何

所谓解析几何，可以称之为懒人的几何

所谓解析几何，可以称之为懒人的几何

所谓解析几何，可以称之为懒人的几何

所谓解析几何，可以称之为懒人的几何

所谓解析几何，可以称之为懒人的几何

所谓解析几何，可以称之为懒人的几何

所谓解析几何，可以称之为懒人的几何

所谓解析几何，可以称之为懒人的几何

所谓解析几何，可以称之为懒人的几何

所谓解析几何，可以称之为懒人的几何

所谓解析几何，可以称之为懒人的几何

所谓解析几何，可以称之为懒人的几何

所谓解析几何，可以称之为懒人的几何

所谓解析几何，可以称之为懒人的几何

所谓解析几何，可以称之为懒人的几何

所谓解析几何，可以称之为懒人的几何

所谓解析几何，可以称之为懒人的几何

所谓解析几何，可以称之为懒人的几何

所谓解析几何，可以称之为懒人的几何

所谓解析几何，可以称之为懒人的几何

所谓解析几何，可以称之为懒人的几何

所谓解析几何，可以称之为懒人的几何

所谓解析几何，可以称之为懒人的几何

所谓解析几何，可以称之为懒人的几何

所谓解析几何，可以称之为懒人的几何

所谓解析几何，可以称之为懒人的几何

所谓解析几何，可以称之为懒人的几何

所谓解析几何，可以称之为懒人的几何

所谓解析几何，可以称之为懒人的几何

所谓解析几何，可以称之为懒人的几何

所谓解析几何，可以称之为懒人的几何

所谓解析几何，可以称之为懒人的几何

所谓解析几何，可以称之为懒人的几何

所谓解析几何，可以称之为懒人的几何

所谓解析几何，可以称之为懒人的几何

所谓解析几何，可以称之为懒人的几何

所谓解析几何，可以称之为懒人的几何

所谓解析几何，可以称之为懒人的几何

所谓解析几何，可以称之为懒人的几何

所谓解析几何，可以称之为懒人的几何

所谓解析几何，可以称之为懒人的几何

所谓解析几何，可以称之为懒人的几何

所谓解析几何，可以称之为懒人的几何

所谓解析几何，可以称之为懒人的几何

所谓解析几何，可以称之为懒人的几何


我们在平面几何和立体几何里, 所用的研究方法是以公理为基础, 直接依据图形的点、线、面的关系来研究 图形的性质. 在将要学习的平面解析几何里, 所用的研究方法和平面几何、立体几何不同, 它是在坐标系的基础 上, 用坐标表示点, 用方程表示曲线 (包括直线), 通过研究方程的特征间接地来研究曲线的性质. 因此可以说, 解 析几何是用代数方法来研究几何问题的一门数学学科.
平面解析几何研究的主要问题是: (1) 根据已知条件, 求出表示平面曲线的方程; (2) 通过方程, 研究平面曲线的性质.
解析几何的这种研究方法, 在进一步学习数学、物理和其他科学技术中经常使用. 在十七世纪, 法国数学家笛卡儿创始了解析几何. 解析几何的产生对数学发展, 特别是对微积分的出现起了 促进作用; 恩格斯对笛卡儿的这一发现给予了高度的评价.
\chapter{解析几何（一）}
\section{平面解析几何简介}
\subsection{平面直角坐标系中的基本公式}
所谓解析几何，就是用坐标研究点，用方程研究曲线。在初中的时候，就已经学习过了数轴，且认定这样的事实：
实数和数轴上的点是一一对应的。那么用实数研究数轴上的点，衡量线段的长度就是自然而然的了。
\vspace*{-5pt}
\begin{figure}[h]
    \centering
    \begin{tikzpicture}
        % Draw the number line
        \draw[->] (-1,0) -- (5,0) node[right] {$x$};
    
        % Draw the points x1 and x2
        \filldraw (1,0) circle (1pt) node[below] {$x_1$} node[above]{$A$};
        \filldraw (3,0) circle (1pt) node[below] {$x_2$} node[above]{$B$};
    \end{tikzpicture}
\end{figure}
\vspace*{-5pt}
\begin{theorem}\label{数轴上的距离}在数轴上有两点\(A,B\)，坐标分别为\(x_1,x_2\)。记\(A,B\)之间的距离为\(d\)，则
    \begin{equation}
        d=|x_1-x_2|
    \end{equation}
\end{theorem}

关于\(A,B\)两点之间的距离，使用\(|AB|,|\ve{AB}|\)都是可以的；而\(AB\)这种表示既可以
表示\(A,B\)两点之间的距离，也可以表示过\(A,B\)两点的直线，使用的时候需要注意避免歧义。在写出
如\(AB\pxqdy CD\)，意味着\(|AB|=|CD|\)，且过\(AB\)两点的直线平行于过\(CD\)两点的直线。

\begin{definition}
    在定理\ref{数轴上的距离}的基础上，再加入一个点\(C\)，使得\(|AC|,|BC|\)之间总保持着一定的比例关系
，则称点\(C\)为定比分点。这个定义可以任意地推广到其他有距离的情形。若\(C\)在\(A,B\)之间，称\(C\)为内分点；
若\(C\)在\(A,B\)两侧，称\(C\)为外分点。
\end{definition}
\begin{theorem}
    在数轴上有\(A,B,C\)三点，坐标分别为\(x_1,x_2,x_3\)。若满足\(\ve{AC}=\lambda\ve{CB} \ (\lambda\neq 1)\)，
    则\begin{equation}
        x_3=\frac{x_{1}+\lambda x_{2}}{1+\lambda} \quad(\lambda \neq-1)
    \end{equation}
\end{theorem}

而在平面上也可以研究相似的问题。在下文中如果没有特殊说明，平面上的坐标总是平面直角坐标系中的坐标，
即在平面内取一组正交单位基底得到的。

\begin{theorem}[两点之间距离公式]
    在平面直角坐标系中，有\(A,B\)两点，坐标分别为\((x_1,y_1),(x_2,y_2)\)。记\(A,B\)之间的距离为\(d\)，则
    \begin{equation}
        d=\sqrt{(x_1-x_2)^2+(y_1-y_2)^2}
    \end{equation}
\end{theorem}
\begin{proof}
    
\end{proof}
\begin{exercise}
    已知三点\(A(3,2),B(0,5),C(4,6)\)，判断\(\tri ABC\)形状。
\end{exercise}



\begin{theorem}
    在数轴上有\(A,B,C\)三点，坐标分别为\((x_1,y_1),(x_2,y_2),(x_3,y_3)\)。若满足\(\ve{AC}=\lambda\ve{CB} \ (\lambda\neq 1)\)，则
    \begin{equation}
        x_3=\frac{x_{1}+\lambda x_{2}}{1+\lambda}, y_3=\frac{y_{1}+\lambda y_{2}}{1+\lambda}
    \end{equation}
\end{theorem}
\begin{corollary}\label{定比分点公式定理}
    点\(P\)为线段\(AB\)的中点，\(A,B\)坐标为$(x_1,y_1),(x_2,y_2)$，则点\(P\)坐标为
    \begin{equation}\label{定比分点公式}
        \left(\frac{x_{1}+x_{2}}{2}, \frac{y_{1}+y_{2}}{2}\right)
    \end{equation}
\end{corollary}
\begin{example}\label{三等分点}
    点\(P\)为线段\(AB\)的三等分点中更靠近\(A\)点的一个，\(A,B\)坐标为$(x_1,y_1),(x_2,y_2)$，
    求\(P\)点坐标。
\end{example}

需要使用定比分点坐标公式的时候，往往不会直接给出形如\(\ve{AC}=\lambda\ve{CB} \ (\lambda\neq 1)\)的式子。
可以根据自己的需求灵活地重新选定起点、终点和分点，
一但选定，\(\lambda\)也就随之而定。因此解题时，必须说明\(\lambda\)的意义。
在解答题中，不要直接写如式\ref{定比分点公式}，而是仿照定理\ref{定比分点公式定理}和例\ref{三等分点}推导想要的形式即可。

\begin{theorem}
    在\(\tri ABC\)中，点\(A,B,C\)的坐标分别为\((x_1,y_1),(x_2,y_2),(x_3,y_3)\)。则\(\tri ABC\)的重心
    坐标为
    \begin{equation}
        \left(\frac{x_{1}+x_{2}+x_{3}}{3},  y=\frac{y_{1}+y_{2}+y_{3}}{3} .\right)
    \end{equation}
\end{theorem}

在学习了平面直角坐标系中的基础——坐标与距离的求法之后，来了解一些简单的应用——
建立适当的坐标系解决几何问题。实际上，上面坐标与距离的知识，以及用坐标解决问题在向量
已经了解很多了，这里只是再拿出来汇总一下。
%坐标法解决问题在向量里将很多了吧，就别在这炒冷饭了，，，
\begin{example}
在$ \parallelogram A B C D $中，求证： 
\begin{equation}
    A C^{2}+B D^{2}=2\left(A B^{2}+A D^{2}\right)
\end{equation}
\end{example}
\begin{exercise}
    已知  $A B C D $ 是一个长方形, 且 $ M $ 是 $ A B C D  $所在平面上任意一个点, 求证: $ A M^{2}+C M^{2}=B M^{2}+D M^{2} $
\end{exercise}


接下来谈一谈“距离”这个词。提起距离，大家总会默认这样的事实：两点之间用线段连接起来的那段长度
才叫距离。但是在现实生活中，总是能走直线吗？在高楼林立的城市中，你必须沿着纵横的街道走（暂时忽视地球是个球这件事）；
在中国象棋的棋盘上，除了马、象，其他棋子总是只能走横纵的直线的。将这种现象抽象成数学语言，就是在平面直角坐标系中有两点\((x_1,y_1),(x_2,y_2)\)，则
\(|x_1-x_2|+|y_1-y_2|\)也一定程度上可以衡量两点之间的距离。

但是还有一个问题等待着回答：究竟什么是距离？依赖于直觉，应该可以给出这样的一些特征：
距离至少不能是负数；距离之间应该是对称的，即从\(A\)到\(B\)的距离应该等于
从\(B\)到\(A\)的距离……
\begin{definition}[距离空间]
    设$X$为非空集合，在$X$上定义一个双变量的实值函数\(\rho:X\times X\rightarrow\mathbb{R}\)，
    若对于任意的\(x,y,z\in X\)满足下列条件：
    \begin{enumerate}
        \setlength{\itemsep}{-\itemsep}
        \item 正定性\(\rho(x,y)\leq 0,\rho(x,y)=0\Leftrightarrow x=y\)
        \item 对称性\(\rho(x,y)=\rho(y,x)\)
        \item 三角不等式\(\rho(x,y)\leq \rho(x,z)+\rho(z,y)\)
    \end{enumerate}
    则称\((X,\rho)\)为一个以\(\rho\) 为距离（或度量）的距离（度量）空间，函数\(\rho\)为其上的一个距离函数。
\end{definition}
\begin{example}求证
  \[ d_1((x_1,y_1),(x_2,y_2))=|x_1-x_2|+|y_1-y_2|\]和
    \[d_2((x_1,y_1),(x_2,y_2))=\sqrt{(x_1-x_2)^2+(y_1-y_2)^2}\]都是\(\mathbb{R}^2\)上的距离（度量）函数。
\end{example}
\begin{proof}
    正定性和对称性都是显然的，仅证明满足三角不等式即可。前者依赖于绝对值不等式，后者由几何意义即证。
    \begin{align*}
        &d_1((x_1,y_1),(x_3,y_3))+ d_1((x_2,y_2),(x_3,y_3))\\
        =&|x_1-x_3|+|y_1-y_3|+|x_3-x_2|+|y_3-y_2|\\
        \leq&|x_1-x_3+x_3-x_2|+|y_1-y_3+y_3-y_2|\\
        =&d_1((x_1,y_1),(x_2,y_2))\\
    \end{align*}
\end{proof}
实际上有这样较强的定理：
\begin{theorem}
    \((\mathbb{R}^n,\rho_p),p\leq 1\)是距离空间，其中
    \[\rho_p(x,y)=\left(\sum_{k=1}^{n}|x_k-y_k|^p\right)^{\frac{1}{p}}\]
    三角不等式由Minkowski不等式保证。
\end{theorem}

当然，这些东西的用处还太过遥远，距离是用来刻画收敛的一个强力工具，在分析学中你还会再见到它，如果不幸学了数学。

\clearpage








\subsection{曲线与方程}
在初中和高一的学习中，已经接触到了许多关于曲线与方程的知识。坐标系沟通了
代表几何的曲线与代表代数的方程。通过几何图形，可以快速认识到（或是简单地猜测）
对称性、取值范围、最值与极值等性质，同时可以反哺代数，有意识地处理一个陌生的方程。

举个具体的例子：还记得你是怎么研究一次函数\(y=2x\)的吗？或许你列了这样的一个表：
\begin{table}[h]\centering
        \begin{tabular}{|c|l|l|l|l|l|}
        \hline
        $x $   &$ -2$ & $-1$ & $0$ &$ 1$ & $2$ \\ \hline
        $y=2x$ & $-4 $&$ -2 $& $0$ & $2$ & $4$ \\ \hline
        \end{tabular}
\end{table}\\然后把它们绘制在了平面直角坐标系中，并用平滑的曲线连接起来——这里是一条直线。虽然你可能不知为什么
要这么做，但是你也大胆的猜测：如果把所有\(x\)可能的取值对应的点都绘制出来，那么一定是这条直线。
这样的研究方法在暗示着一件事：符合方程的坐标对应的点在方程对应的曲线上，完成了从代数到几何的跨越。
那你立刻就可以说出另外一个方向——从几何到代数：在曲线上的点对应的坐标符合方程。由此，抽象出下面的定义：

\begin{definition}\label{方程的曲线}
    一般地，在平面直角坐标系中，如果某曲线$C$（看作点的集合或适合某种条件的点的轨迹）上的点与一个二元
方程\[f(x,y)=0\]的实数解建立了如下的关系：
\vspace*{-8pt}
\begin{enumerate}
    %\setlength{\itemsep}{0pt}
    \setlength{\itemsep}{-\itemsep}
    \item 曲线上点的坐标都是这个方程的解；
    \item 以这个方程的解为坐标的点都是曲线上的点．
\end{enumerate}
\vspace*{-8pt}
那么,这个方程叫做曲线的方程；这条曲线叫做方程的曲线(curve)，有时也称方程的轨迹。

此时可以将曲线$C$看做是如下的一个集合：
\vspace*{-10pt}
\[\{(x,y)\in\mathbb{R}^2|f(x,y)=0\}\]
\end{definition}
利用这个定义，可以解决下面的例题与练习。
\begin{example}判断下列说法正误，并说出理由。

    $(1)$三角形的一条中线上至少有两个点在$x$轴上，所以这条中线的方程是$y=0$。

    \((2)\)曲线\(C\)上无数个点的坐标都满足方程$F(x,y)=0$，
    方程$F(x,y)=0$的所有实数解为坐标的点都在曲线$C$
    上，所以曲线\(C\)是方程$F(x,y)=0$的曲线。
\end{example}
\begin{jie}
    两句话都不对。

    \((1)\)三角线中线是一条线段，中线上两点在\(x\)轴上，只说明中线在\(x\)轴上，或中线所在的直线是\(x\)轴，但中线不是\(x\)轴。
    换句话说，在三角形外的轴上的点，它的坐标为\((x_0,0)\)，满足方程$y=0$，但它不
    在中线上，不满足曲线与方程定义的第二个条件，所以不能说中线的方程为\(y=0\)。

    \((2)\) 曲线 $ C $ 上无数点的坐标都满足方程$  F(x, y)   =0$ , 
    不能保证曲线  $C $ 上每一个点的坐标都满足方程  $F(x ,  y)=0$ 。
    用\((1)\)的例子即可，考虑\(y=0,\ x\in(-1,1)\)作为\(F(x,y)\)，\(x\)轴作为
    曲线\(C\)。因此，不满足曲线与方程定义的第一个条件， 不能说曲线  $C$  是方程 $ F(x, y)=0$  的曲线。
\qed
\end{jie}

通过这个例子是想要说明，定义\ref{方程的曲线}并不是废话。从集合的观点，曲线可以看作是具有
某种本质特征（条件）的点的集合，或者看作是动点按照某种运动规律运动的轨迹，但是总要通过严格的定义保证：
这条曲线上的点都具有这种本质特征 \textcolor{red}{（纯粹性）}，具有这种本质特征的点都在这条曲线上 \textcolor{red}{（完备性）}。
\begin{example}
若\((x,2)\)在曲线\(x^2+xy+y^3=9\)上，求\(x\)的值。
\end{example}
\begin{jie}
    点在曲线上，故点对应的坐标是方程的解。带入其中得到
    \[x^2+2y+8=9\]

    解得\(x=-1\pm\sqrt{2}\)
    \qed
\end{jie}

\begin{example}
    求方程\((x-y)^2+(xy-1)^2=0\)表示的曲线。
\end{example}
\begin{jie}
    由于\((x-y)^2\leq 0,(xy-1)^2\leq 0\)，又\((x-y)^2+(xy-1)^2=0\)，故总有
    \((x-y)^2=0\)且\((xy-1)^2=0\)。仅有\((1,1),(-1,-1)\)符合。
    \qed
\end{jie}
\begin{exercise}
    若曲线\(ax^3+by^2=4\)经过了\((0,2)\)和\((\sqrt[3]{2},1)\)，求该曲线方程。
\end{exercise}
\begin{exercise}
    求方程\(x^2+4xy+4y^2=9\)表示的曲线。
\end{exercise}

接下来了解一下，如果根据曲线求出曲线对应的方程。

\begin{example}
    求平分一、三象限的直线方程。
\end{example}
\begin{jie}
    由角平分线的判定定理，在角的内部到角的两边距离相等的点在角平分线上。设角平分线上点为\(P(x,y)\)，故
    这些点构成的集合是\[\{P|d_{P-x\text{轴}}=d_{P-y\text{轴}}\}\]
    点\((x,y)\)到\(x\)轴
    距离为\(|y|\)，到\(y\)轴距离为\(|x|\).故
    \[|x|=|y|\]
    解出来\(x=\pm y\)，其中\(y=-x\)不在一三象限的内部，舍去。仅考虑\(y=x\)：符合方程的坐标对应的点在一、三角平分线上，
    一、三象限角平分线上的点也符合\(y=x\)。故平分一、三象限的直线方程为\(y=x\)。
    \qed
\end{jie}
\begin{comment}
解出来两条直线，一条是内角平分线，另一条是外角平分线。在这里\(y=-x\)是第二、四象限角平分线。
\end{comment}
归纳出如下步骤：
\tbox{
\begin{enumerate}
    \setlength{\itemsep}{-\itemsep}
    \item 建立适当的坐标系，用有序实数对\((x,y)\)表示曲线上任意一点M的坐标；
    \item 写出适合条件\(p\)的点\(M\)的集合 $P=\{M|p(M)\}$；
    \item 用坐标表示条件$p(M)$，列出方程\(f(x,y)=0\)；
    \item 化方程\(f(x,y)=0\)为最简形式；
    \item 说明以化简后的方程的解为坐标的点都在曲线上，如果某个点是方程的解却不在曲线上要剔除。
\end{enumerate}
    }
\begin{example}
    已知两点 $ A(-2,-2), B(2,2)$ ，试求满足条  $||M A|-|M B||=4$  的动点  M  的轨迹方程。
\end{example}
\begin{jie}
    设动点  $M $ 的坐标为 $ (x, y)$  ，则由条件$||M A|-|M B||=4$
    可得
    
   $$ |\sqrt{(x+2)^{2}+(y+2)^{2}}-\sqrt{(x-2)^{2}+(y-2)^{2}}|=4,$$
    移项得
    
   $$ \sqrt{(x+2)^{2}+(y+2)^{2}}=\sqrt{(x-2)^{2}+(y-2)^{2}}\pm 4$$
    两边平方, 整理得
    
  $$\pm \sqrt{(x-2)^{2}+(y-2)^{2}}=x+y-2 $$
    再次平方得
   $$ x y=2$$

   其上一点\(P\)不妨表示为\((x,\dfrac{2}{x})\)，到\(A,B\)两点距离为
   \[|PA|=\sqrt{(x+2)^2+(\frac{2}{x}+2)^2}=\sqrt{(x+\frac{2}{x}+2)^2}=|x+\frac{2}{x}+2|\]
    \[|PB|=\sqrt{(x-2)^2+(\frac{2}{x}-2)^2}=\sqrt{(x+\frac{2}{x}-2)^2}=|x+\frac{2}{x}-2|\]
    由基本不等式可知，\(x+\dfrac{2}{x}\geq 2\sqrt{2}\)或\(x+\dfrac{2}{x}\leq -2\sqrt{2}\)，由此脱下绝对值，总满足
    \[||M A|-|M B||=4\]
    故以方程的解为坐标的点都满足该动点的条件，\(xy=2\)就是所求轨迹的方程。
    \qed
\end{jie}
方程化简到什么程度，并没有明确的规定，一般指
将方程\(f(x,y)=0\)化成关于\(x,y\)的整式。如果化简过程破坏了同解
性，就需要剔除不属于轨迹上的点，找回属于轨迹而遗漏的点。化简方程时候，要尽量保证方程的同解性，在平方、去绝对值、去分母的时候要尤其
谨慎。

利用下面两个练习，初步学习如何将几何条件翻译为坐标语言。
\begin{exercise}
    在平面直角坐标系中，\(A\ (-1,-1),B\ (3,7)\)，求线段\(AB\)的垂直平分线的方程。
\end{exercise}
\begin{hint}
    到线段两端距离相等的点在线段的垂直平分线上。
\end{hint}
\begin{exercise}
    在平面直角坐标系中，已知点$M(4,0),N(1,0)$，若动点P满足$\overrightarrow{MN}\cdot\overrightarrow{MP}=6|\overrightarrow{NP}|$。求
动点$P$的轨迹$C$的方程。
\end{exercise}
\begin{hint}
    向量关系可以翻译为坐标关系。
\end{hint}
\begin{exercise}
    下列方程表示同一曲线的一组是
    \begin{enumerate}
        \setlength{\itemsep}{-\itemsep}
        \item \(xy=1\)，\(y=\dfrac{1}{x}\)
        \item \(y=|x|\)，\(y=\pm x\)
        \item \(y=\dfrac{x^2-1}{x-1}\)，\(y=x+1\)
        \item \(y=\sqrt{1-x^2}\)，\(x^2+y^2=1\)
    \end{enumerate}
\end{exercise}

接下来在认识一下，如何借助已知轨迹方程，去求其他的轨迹方程。
\begin{example}在平面直角坐标系中，\(O\)为坐标原点。已知圆$C:x^2+(y-3)^2=9$，过原点作圆$C$的弦$OP$，
求$OP$的中点$Q$的轨迹方程。
\end{example}
\begin{jie}

\end{jie}
\begin{exercise}
    动点\(P\)为曲线\(y=2x^2+1\)上任意一点，定点\(A(0,-1)\)，线段\(PA\)上有一点\(M\)，
    满足\(PM:AM=1:2\)，求\(M\)点轨迹方程。
\end{exercise}
\tbox{用相关点法求解轨迹：
\begin{enumerate}
    \setlength{\itemsep}{-\itemsep}
    \item 设点\((x,y)\)待求轨迹上，点\((x_0,y_0)\)在已知轨迹\(f(x,y)=0\)上；
    \item 找到\((x,y)\)与\((x_0,y_0)\)的关系，列出方程；
    \item 从中解出
       $ \begin{cases}
         x_0= g_1(x,y)\\
         y_0= g_2(x,y) 
      \end{cases}$；
    \item 带入\(f(x_0,y_0)=0\)，得到\(f(g_1(x,y),g_2(x,y))=0\)，化简得到关于\(x,y\)的轨迹方程。
\end{enumerate}
}

在这一节的最后，以曲线的交点做结。如果一个点是两条曲线的公共点，
那么点的坐标也必然是两条曲线对应方程的公共解。反过来说，如果想要求得几条曲线的交点，那么只需要联立
对应的方程，求取公共解即可。
\begin{example}
    求直线 $ y=x+\dfrac{3}{2} $ 被曲线  $y=\dfrac{1}{2} x^{2}$  截得的线段的长。
\end{example}
\begin{jie}
    
\end{jie}
\clearpage
\review
\clearpage
\hw
\begin{homework}
    证明与两条坐标轴的距离的积是常数$k(k>0)$的点的轨迹方程是$xy=\pm k$。
\end{homework}
\begin{homework}
    在平面直角坐标系中，等腰三角形\(ABC\)的底边两端点是\(A\ (1,0),B\ (-1,0)\)，顶点\(C\)为一动点，试求C的轨迹。
\end{homework}
\clearpage
